\chapter{การประมวลผล Big Data และโมเดลการเรียนรู้ของเครื่องด้วย Dataproc Serverless}

สำหรับงานประมวลผลข้อมูลขนาดใหญ่ที่ไม่เหมาะกับโครงสร้างตารางเชิงสัมพันธ์ Apache Spark บน Dataproc Serverless มอบพลังการคำนวณแบบกระจายศูนย์โดยไม่ต้องเสียเวลาจัดการคลัสเตอร์เสมือน

\section{สถาปัตยกรรม Apache Spark Serverless}

ผู้ใช้งานเพียงส่งไฟล์สคริปต์ PySpark และกำหนดคุณลักษณะของทรัพยากร ระบบของ Google Cloud จะทำการจัดสรรหน่วยคำนวณและปิดระบบลงอัตโนมัติเมื่อประมวลผลเสร็จสิ้น (Just-in-Time Provisioning) ทำให้คิดค่าใช้จ่ายตามระยะเวลาการทำงานจริงระดับวินาที

\section{การบริหารหน่วยความจำเพื่อป้องกันข้อผิดพลาด Out-Of-Memory}

การจัดสรรหน่วยความจำของ Spark Executor ต้องคำนึงถึงสัดส่วนระหว่างพื้นที่จัดเก็บ (Storage Memory) และพื้นที่คำนวณ (Execution Memory):

\begin{equation}
\text{Total Executor Memory} = \text{Overhead} + \text{Off-Heap} + \text{On-Heap} (\text{Storage} + \text{Execution} + \text{User})
\end{equation}

การใช้คำสั่ง \texttt{repartition()} อย่างเหมาะสม ร่วมกับการหลีกเลี่ยงคำสั่ง \texttt{collect()} บนข้อมูลขนาดใหญ่ ช่วยให้การคำนวณไหลลื่นโดยไม่เกิดอาการ Driver Crash

\section{การสร้างไปป์ไลน์ Machine Learning ด้วย Spark MLlib}

การฝึกสอนโมเดลทำนาย เช่น การพยากรณ์ความชื้นในดินล่วงหน้าด้วย Random Forest Regressor สามารถรันบนคลัสเตอร์หลายร้อยคอร์พร้อมกัน ทำให้ลดเวลาการฝึกสอนโมเดลจากหลายชั่วโมงเหลือเพียงไม่กี่นาที
